\documentclass[../../main.tex]{subfiles}
\begin{document}

{\bf [解]}

\paragraph{(1)}

題意の条件より
\begin{align}
 |\vec a| = |\vec b| = 1 \label{eq:1}
\end{align}
である．
$\vec a,\vec b$のなす角$\theta$とおくと
\begin{align}
 \vec{a}\cdot\vec{b} = \cos\theta \label{eq:2}
\end{align}
である．
対称性から
\begin{align}
 0<\theta\le\dfrac{\pi}{2} \label{eq:3}
\end{align}
で考えれば良い．
というのも，$\dfrac{\pi}{2}<\theta<\pi$の時は$\vec{b'}=-\vec{b},n'=-n$と置き直せば\cref{eq:3}の場合に帰着するからである．
この時\cref{eq:2}で$\cos\theta \ge 0$であることに気をつけよう．

さて，$(m,n)\neq (0,0)$に対して
\begin{align}
 |l(m,n)|^2=m^2+n^2+2mn\cos\theta  \label{eq:4}
\end{align}
の最小値を考えよう．対称性から$m \ge 0$として考えれば十分である．
$m<0$の場合は$l(m,n)=-l(-m,-n)$に注意すれば$m>0$の場合の議論と合流するからである．

$m=0$の時は\cref{eq:4}は
\begin{align}
 |l(0,n)|^2=n^2 \label{eq:5}
\end{align}
であり，これは自明に$n= \pm 1$で最小値$1$をとる．
そこで以下$m>0$の場合について考える．

\subparagraph{$n\ge 0$の時}

この場合は\cref{eq:4}の各項が$m,n$について単調増加だから$m,n$を小さくした時が最小値であり，$(m,n)=(1,0)$で最小値$1$をとる．

\subparagraph{$n<0$の時}

この場合は\cref{eq:4}の第三項が負になる．そこで平方完成すると
\begin{align}
|l(m,n)|^2
 &=\left(n+m\cos\theta\right)^2 + m^2 - m^2\cos^2\theta \\
 &=\left(n+m\cos\theta\right)^2 + m^2\left(1-\cos^2\theta\right) \\
 &=\left(n+m\cos\theta\right)^2 + m^2\left(1-\cos\theta\right)\left(1+\cos\theta\right) \label{eq:6}
\end{align}
である．

ここで$m\ge 2$の時は第二項を下から評価して
\begin{align}
 |l(m,n)|^2
  &\ge \left(n+m\cos\theta\right)^2 + 4\left(1-\cos\theta\right)\left(1+\cos\theta\right) \\
  & >  \left(n+m\cos\theta\right)^2 + 2\left(1-\cos\theta\right) \\
  & =  \left(n+m\cos\theta\right)^2 + \left|\vec{a}-\vec{b}\right|^2 \\
  & \ge \left|\vec{a}-\vec{b}\right|^2 \\
  & = |l(1,-1)|^2 \\
\therefore
 |l(m,n)|^2 > |l(1,-1)|^2 \label{eq:7}
\end{align}
である．

全く同様に$n \le -2$の時は$m$について平方完成すれば
\begin{align}
  |l(m,n)|^2 > |l(1,-1)|^2 \label{eq:8} 
\end{align}
である．従って最小値は残る$(m,n)=(1,-1)$の時に取る．
\casesend

以上の場合わけで全ての場合が尽くされた．よって\cref{eq:3}および$m\ge 0$の場合の最小値は
$(m,n)=(0,\pm 1),(1,0),(1,-1)$のいずれかで取る．

$\dfrac{\pi}{2}<\theta<\pi$の場合に範囲を拡大すると，対称性から$(m',n')=(m,-n)$の場合も考えれば良い．
従ってさらに$(m,n)=(1,1)$の場合も含む．これらのうち$(1,0), (0,\pm 1)$の場合の値は$1$だから，結局
\begin{align}
 r & = \min l(m,n) \\
   & = 1, l(1,-1), l(1,1) \\
   & = 1, |\vec{a}-\vec{b}|, |\vec{a}+\vec{b}| 
\end{align}
であり題意は示された．

\paragraph{(2)}

引き続き\cref{eq:3}および$m\ge 0$のもとで考えれば十分である．
この時$r=1$となる条件は$|\vec{a}-\vec{b}|$がこれより大きいことで
\begin{align}
  & |\vec a-\vec b|\ge1 \\
\iff 
  & 2-2\cos\theta\ge1 \\
\iff 
  & \dfrac{1}{2} \ge \cos\theta \\
\iff
  & \dfrac{\pi}{3}\le\theta\le\dfrac{\pi}{2} \label{eq:9}
\end{align}
となる．従って三角形$OAB$の面積
\begin{align}
 S = \dfrac{1}{2}\sin\theta
\end{align}
の最小値は$\theta=\dfrac{\pi}{3}$の時の$\dfrac{\sqrt{3}}{4}$である．

\newpage

\end{document}
